\section{Информационные состояния}

В работе сравниваются четыре информационных состояния.

\paragraph{Только агрегаты.}
\[
s_t^{agg} = (\pi_t^{obs}, y_t^{obs}).
\]

\paragraph{Оценённые агрегаты.}
\[
s_t^{filter}
=
\mathbb E[
(\pi_t^{gap},y_t^{gap},r_t^*)
\mid \mathcal I_t].
\]

\paragraph{Агрегаты и распределение.}
\[
s_t^{dist}
=
\left(
\hat \pi_t,
\hat y_t,
\hat r_t^*,
\widehat{MPC}_t,
\widehat{share}^{liq}_t
\right).
\]

\paragraph{Полная информация.}
\[
s_t^{full}=x_t.
\]

Полная информация используется только как верхняя граница качества правила.
